\section{Introducción}
El marco competitivo del sector cárnico impone procesos de talento humano ágiles, trazables y accesibles a través del tiempo. En "Las Brisas", se contaba con información tanto en archivos como en planillas teniendo en cuenta que había duplicidades y muy poco control del cambio de información. La digitalización de estos procesos hace posible eliminar reprocesos, acelerando así la emisión de certificados y mejorar el cumplimiento normativo. La literatura revisada apunta a beneficios por la adopción de soluciones digitales en RR. HH.,\cite{analisis_lenguajes} además de los inconvenientes por el proceso de adopción o, en su defecto, cambio cultural \cite{gestion_talento_digital}.

En cuanto al nivel tecnológico, se eligió un stack suficientemente soportado: React en el frontend por su arquitectura basada en componentes y su ecosistema, y Spring Boot en el backend por su madurez, seguridad y productividad en entornos empresariales \cite{analisis_frameworks_frontend}\cite{estudio_aplicaciones_java}. En su interfaz, se optó diseño responsivo de acuerdo con recomendaciones para la accesibilidad y lógica en el uso de dispositivos \cite{optimizacion_bootstrap}. La opción de MySQL se validó a partir principios de modelado relacional y normalización \cite{analisis_google_charts}. Finalmente, se espera que se pueda ir evolucionando hacia una ruta de microservicios para módulos con alta lectura \cite{implementacion_microservicios}.